\documentclass[11pt]{article}
\usepackage[a4paper,margin=1in]{geometry}
\usepackage[T1]{fontenc}
\usepackage[utf8]{inputenc}
\usepackage[ngerman]{babel}
\usepackage{lmodern}
\usepackage{microtype}
\usepackage{amsmath,amssymb,mathtools}
\usepackage{enumitem}
\setlength{\parindent}{1.2em}
\setlength{\parskip}{0.35em}
\setlength{\emergencystretch}{5em}
\allowdisplaybreaks
\newcommand{\dd}{\,\mathrm d}
\newcommand{\ii}{\mathrm i}
\newcommand{\cosec}{\operatorname{cosec}}
\newcommand{\tg}{\operatorname{tg}}
\newcommand{\arccosop}{\operatorname{arc\,cos}}
\newcommand{\art}[1]{\bigskip\begin{center}\large\bfseries #1.\end{center}\nopagebreak}
\begin{document}
\begin{center}
{\Large I.\\[0.4em] Über die Elemente der Theorie der Eulerschen Integrale.}\par
\medskip
{\normalsize Fortsetzung: Artikel 9--12.}
\end{center}

\art{9}

Ein von den bisher angeführten wesentlich verschiedener Beweis ist endlich noch in der Abhandlung „Disquisitiones generales circa seriem infinitam etc. Auctore C. F. Gauß“ enthalten. Die ganze Anlage derselben macht es aber unmöglich, diesen Beweis hier darzustellen, indem die Grundlagen, auf welche er sich stützt, mit einem Schlage eine vollständige Theorie der Eulerschen Integrale ergeben, deshalb aber auch zu bedeutend sind, um hier bloß zu diesem einzigen Zwecke entwickelt zu werden.

Zu diesen will ich nun noch einen, so viel mir bekannt ist, neuen Beweis hinzufügen, der eben nicht viel Zurüstungen erfordert und sich stets in dem Gebiete der Integralrechnung hält, wenn auch die Methode nicht eben neu ist, die darauf hinaus läuft, die Aufgabe auf die Integration einer Differentialgleichung zweiter Ordnung zurückzuführen. Setzt man in der Gleichung
\[
B B=\int_0^\infty\frac{x^{b-1}\dd x}{x+1}
     \int_0^\infty\frac{y^{b-1}\dd y}{y+1}
=\int_0^\infty\frac{\dd x}{x+1}
     \int_0^\infty\frac{(xy)^{b-1}}{y+1}\dd y
\]
$y=\frac zx$, $\dd y=\frac{\dd z}{x}$, kehrt dann die Integrationsordnung um, und führt die Integration in bezug auf $x$ aus, so findet man leicht
\begin{equation}
B B=\int_0^\infty\frac{z^{b-1}\ell z}{z-1}\dd z
=\frac{\dd}{\dd b}\int_0^\infty\frac{z^{b-1}\dd z}{z-1}.
\tag{27}
\end{equation}
Durch unbestimmte Integration in bezug auf $b$ erhält man daher
\[
\int B B\,\dd b=\int_0^\infty\frac{z^{b-1}\dd z}{z-1},
\]
worin das Integral rechts sich bloß durch die Form des Nenners von dem Integral $B$ unterscheidet; und es ist leicht vorauszusehen, daß man das Integral $B$ wiedererhält, wenn man das eben gewonnene derselben Behandlung unterwirft. In der Tat findet man
\[
B\int B B\,\dd b
=\int_0^\infty\frac{\dd z}{z-1}
  \int_0^\infty\frac{(zy)^{b-1}}{y+1}\dd y,
\]
und wenn man $y=\frac xz$, $\dd y=\frac{\dd x}{z}$ setzt, dann zufolge Art. 6 und 7 das in bezug auf $z$ genommene Integral in zwei Integrale zerlegt, deren Grenzen $0,1-\delta$ und $1+\varepsilon,\infty$ sind, die Integrationsordnung umkehrt, und die Integration in bezug auf $z$ ausführt, so erhält man
\[
B\int B B\,\dd b
=\int_0^\infty\frac{x^{b-1}\ell x}{x+1}\dd x
+\lim\int_0^\infty\frac{x^{b-1}\dd x}{x+1}\,\ell\!\left(\frac{\delta}{\varepsilon}\right).
\]
Hierin ist nun zwar $\lim \ell(\delta/\varepsilon)$ unbestimmt, jedenfalls aber unabhängig von $x$ und $b$, und mag mit $k$ bezeichnet werden. Dann gibt die letzte Gleichung
\[
B\int B B\,\dd b=\frac{\dd B}{\dd b}+kB
\]
oder, wenn man bedenkt, daß in dem Integral links doch schon eine willkürliche Konstante enthalten ist,
\begin{equation}
B\int B B\,\dd b=\frac{\dd B}{\dd b},
\tag{28}
\end{equation}
woraus man durch Division mit $B$ und Differentiation in bezug auf $b$ die Differentialgleichung zweiter Ordnung
\[
B B=\frac1B\frac{\dd\dd B}{\dd b^2}
-\frac1{B B}\left(\frac{\dd B}{\dd b}\right)^2
\]
erhält, welche die Eigentümlichkeit besitzt, daß die unabhängige Variable $b$ nicht vorkommt, und sich deshalb bekanntlich auf eine Differentialgleichung erster Ordnung zurückführen läßt, wenn man das erste Differentialverhältnis als neue Variable einführt. Bezeichnen wir dieses mit $B'$, so ist
\[
\frac{\dd B}{\dd b}=B',\qquad
\frac{\dd\dd B}{\dd b^2}=\frac{\dd B'}{\dd b}
=\frac{B'\dd B'}{\dd B},
\]
und führen wir diese Transformationen in die obige Differentialgleichung ein, so geht diese in
\[
B B=\frac{B'\dd B'}{B\,\dd B}-\frac{B'B'}{B B}
\]
oder in
\[
\dd(BB)=
\frac{BB\,\dd(B'B')-B'B'\dd(BB)}{(BB)^2}
=\dd\frac{B'B'}{BB}
\]
über, deren Integral
\[
BB=cc+\frac{B'B'}{BB}
=cc+\frac1{BB}\left(\frac{\dd B}{\dd b}\right)^2
\]
ist. Zu der Bestimmung der Konstanten ist nun die Kenntnis zweier Eigenschaften des Integrals erforderlich; diese nehmen wir aus Art. 3, wo gezeigt ist, daß $B$ für $b=\frac12$ ein Minimum $=\pi$ erreicht; da gleichzeitig $\frac{\dd B}{\dd b}=0$ wird, so ergibt sich unmittelbar $cc=\pi\pi$.

Man erhält dann weiter
\[
\pm\,\dd b=
\frac{\dd B}{B\sqrt{BB-\pi\pi}}
=\frac1\pi
  \frac{-\dd\frac{\pi}{B}}
       {\sqrt{1-\frac{\pi\pi}{BB}}},
\]
folglich durch Integration
\[
\pm(b+c)=\frac1\pi\,\arccosop\frac{\pi}{B},
\qquad
B=\frac{\pi}{\cos(b+c)\pi}.
\]
Da $B$ für $b=\frac12$ den Wert $\pi$ erhält, so muß
\[
\cos\left(\frac12+c\right)\pi=-\sin c\pi=1,
\qquad \cos c\pi=0
\]
sein; folglich ist
\[
\cos(b+c)\pi=\cos b\pi\cos c\pi-\sin b\pi\sin c\pi=\sin b\pi,
\]
wodurch man wieder $B=\frac{\pi}{\sin b\pi}$ findet.

\art{10}

Weil der im vorigen Artikel gegebene Beweis den Anforderungen der größten Strenge doch noch nicht Genüge leistet, namentlich insofern das Integral
\[
\int_0^\infty\frac{x^{b-1}\dd x}{x-1}
\]
darin eine Rolle spielt, so ist es wohl nicht unangemessen, hier eine solche Modifikation noch folgen zu lassen, in welcher die unbestimmte Integration in bezug auf $b$ ganz vermieden wird. Da die Gleichung (28) sich auch so schreiben läßt
\[
\int B B\,\dd b=\frac{\dd\,\ell B}{\dd b},
\]
so läßt sich vermuten, daß eine Entwicklung des rechts stehenden Ausdrucks ebenfalls zum Ziele führt. Da $B=\Gamma(b)\Gamma(1-b)$ ist, so reicht es hin, ein Integral für $\frac{\dd\,\ell\Gamma(\mu)}{\dd\mu}$ zu finden, was sich bekanntlich auf folgendem Wege erreichen läßt. Aus der Definition von $\Gamma(\mu)$ folgt
\[
\frac{\dd\,\ell\Gamma(\mu)}{\dd\mu}
=\frac1{\Gamma(\mu)}
  \int_0^\infty x^{\mu-1}e^{-x}\ell x\,\dd x.
\]
Setzt man hierin für $\ell x$ das Integral
\[
\ell x=\int_0^\infty\frac{e^{-z}-e^{-zx}}{z}\dd z,
\]
welches sich aus dem Integral $\int_0^1y^{x-1}\dd y=\frac1x$ durch Integration in bezug auf $x$ zwischen den Grenzen $1$ und $x$, und durch Substitution von $e^{-z}$ für $y$ ergibt, so findet man durch Umkehrung der Integrationsordnung und mit Hilfe von Gleichung (3) sehr leicht
\[
\frac{\dd\,\ell\Gamma(\mu)}{\dd\mu}
=\int_0^\infty\frac{\dd z}{z}
\left(e^{-z}-\frac1{(z+1)^\mu}\right).
\]
Setzt man hierin $b$ und $(1-b)$ für $\mu$, so ergibt sich
\[
\frac{\dd\,\ell B}{\dd b}
=\int_0^\infty\frac{\dd z}{z}
\left(\frac1{(z+1)^{1-b}}-\frac1{(z+1)^b}\right),
\]
und wenn man hierin $z=x-1$ oder $z=\frac1x-1$ setzt:
\[
\frac{\dd\,\ell B}{\dd b}
=\int_1^\infty\frac{x^b-x^{1-b}}{x-1}\frac{\dd x}{x}
=\int_0^1\frac{x^b-x^{1-b}}{x-1}\frac{\dd x}{x},
\]
also auch
\[
2\frac{\dd\,\ell B}{\dd b}
=\int_0^\infty\frac{x^b-x^{1-b}}{x-1}\frac{\dd x}{x}.
\]
Vergleicht man dies mit Gleichung (27), so ergibt sich augenblicklich
\begin{equation}
2\frac{\dd\,\ell B}{\dd b}
=\int_{1-b}^{b} B B\,\dd b
=2\int_{\frac12}^{b} B B\,\dd b,
\tag{29}
\end{equation}
indem ja zufolge Art. 3 $B$ für $b=\frac12+b'$ und $b=\frac12-b'$ dieselben Werte erhält; durch Differentiation dieser Gleichung in bezug auf $b$ erhält man wieder die im vorigen Artikel behandelte Differentialgleichung.

\art{11}

Nachdem durch die angegebenen Beweise die Richtigkeit der Gleichung
\[
B(r,1-r)=\int_0^\infty\frac{x^{r-1}\dd x}{x+1}
=\frac{\pi}{\sin r\pi},
\]
worin $r$ einen positiven echten Bruch bezeichnet, außer Zweifel gesetzt ist, ergibt sich unmittelbar aus Art. 2, daß die Eulerschen Integrale der ersten Art, deren beide Argumente eine ganze positive Zahl zur Summe haben, sich wirklich darstellen lassen; man findet
\begin{equation}
B(m+r,n-r)=
\frac{(m+r-1)\cdots(1+r)r\,(n-1-r)\cdots(2-r)(1-r)}
     {(m+n-1)(m+n-2)\cdots2\cdot1}
\frac{\pi}{\sin r\pi}.
\tag{30}
\end{equation}
Man kann hiervon auch noch einen wichtigen Rückschluß auf die Eulerschen Integrale der ersten Art machen; setzt man nämlich $m=0$, $n=1$, $r=\frac12$, so findet man
\begin{equation}
\Gamma\!\left(\frac12\right)
=\int_0^\infty e^{-x}\frac{\dd x}{\sqrt{x}}
=2\int_0^\infty e^{-xx}\dd x=\sqrt{\pi}
\tag{31}
\end{equation}
und folglich nach Gleichung (4):
\begin{equation}
\Gamma\!\left(n+\frac12\right)
=(2n-1)(2n-3)\cdots5\cdot3\cdot1\,\frac{\sqrt{\pi}}{2^n}.
\tag{32}
\end{equation}

\art{12}

Nachdem in Artt. 2 und 11 die Fälle zusammengestellt sind, in welchen die Eulerschen Integrale beider Arten ohne Hilfe neuer Funktionen dargestellt werden können, will ich zum Schluß noch einmal zu dem Integral $B$ zurückkehren, um noch einige Beziehungen desselben zu anderen Integralen zu entwickeln. Unter den verschiedenen Formen, in welchen es auftritt, sind die beiden folgenden
\[
\int_{-\infty}^{+\infty}
\frac{e^{(2b-1)z}+e^{-(2b-1)z}}{e^z+e^{-z}}\dd z
\quad\text{und}\quad
2\int_0^{\frac{\pi}{2}}(\tg\varphi)^{2b-1}\dd\varphi
=2\int_0^{\frac{\pi}{2}}(\tg\varphi)^{1-2b}\dd\varphi
\]
ganz interessant; wichtiger sind aber die Verallgemeinerungen desselben durch Einführung neuer Konstanten. Dahin gehört das Integral
\begin{equation}
\int_0^\infty\frac{x^{b-1}\dd x}{x+c}
=\frac{\pi c^{b-1}}{\sin b\pi},
\tag{33}
\end{equation}
worin $c$ positiv sein muß; doch läßt sich nachträglich beweisen, daß $c$ auch imaginär sein darf; multipliziert man nämlich Zähler und Nenner der Funktion $\frac{x^{b-1}}{x+i}$ mit $(x-i)$, so zerfällt das entsprechende Integral in zwei andere, welche sich durch die Substitution $xx=y$ auf das Integral $B$ reduzieren lassen, und so findet man die Gültigkeit der Gleichung (33) für imaginäre $c$. Durch Differentiation und Integration in bezug auf $c$ kann man dann wieder eine Reihe von anderen Integralen ableiten.

Sind $c_1,c_2,\ldots,c_n$ voneinander verschiedene positive oder imaginäre Konstanten, ferner $0<b<n$ und
\[
f(x)=(x+c_1)(x+c_2)\cdots(x+c_n),\qquad
f'(x)=\frac{\dd f(x)}{\dd x},
\]
so findet man auch
\begin{equation}
\int_0^\infty\frac{x^{b-1}\dd x}{f(x)}
=\frac{\pi}{\sin b\pi}
 \sum_{k=1}^{n}\frac{c_k^{\,b-1}}{f'(-c_k)},
\tag{34}
\end{equation}
indem man $\frac{x^\beta}{f(x)}$ in Partialbrüche zerlegt, worin $\beta$ die größte in $b$ enthaltene ganze Zahl bedeutet.

Ein Vorzug des in Art. 10 gegebenen Beweises besteht auch noch darin, daß er unmittelbar eine verwandte Klasse von Integralen bestimmen lehrt. Aus den Gleichungen (27) und (29) folgt nämlich unmittelbar
\[
\int_0^\infty\frac{x^b-x^{\frac12}}{1-x}\frac{\dd x}{x}
=\pi\cot b\pi,
\]
folglich auch
\begin{equation}
\int_0^\infty\frac{x^a-x^b}{1-x}\frac{\dd x}{x}
=\pi(\cot a\pi-\cot b\pi),
\tag{35}
\end{equation}
und hieraus ergibt sich auch das in Art. 3 mit $\varphi(b)$ bezeichnete Integral
\[
\begin{aligned}
\int_0^\infty\frac{1-x^\mu}{1-x^\nu}x^{b-1}\dd x
&=\frac1\nu\int_0^\infty
  \frac{x^{\frac b\nu}-x^{\frac{\mu+b}{\nu}}}{1-x}\frac{\dd x}{x}\\
&=\frac{\pi}{\nu}
  \frac{\sin\frac{\mu\pi}{\nu}}
       {\sin\frac{b\pi}{\nu}\sin\frac{(\mu+b)\pi}{\nu}}\\
&=\frac{2\pi}{\nu}
  \frac{\sin\frac{\mu\pi}{\nu}}
       {\cos\frac{\mu\pi}{\nu}-\cos\frac{(\mu+2b)\pi}{\nu}},
\end{aligned}
\]
und das Minimum desselben
\[
\int_0^\infty
\frac{x^{\frac{\mu}{2}}-x^{-\frac{\mu}{2}}}
     {x^{\frac{\nu}{2}}-x^{-\frac{\nu}{2}}}
\frac{\dd x}{x}
=\frac{2\pi}{\nu}\,\tg\frac{\mu\pi}{2\nu}.
\]

Die meisten der hierher gehörigen Integrale finden sich in der im Auftrage des preußischen Ministeriums von Minding herausgegebenen „Sammlung von Integraltafeln“ (Berlin 1849), im Anfang der vierten Abteilung. Vielleicht ist hier der Ort, um einige Fehler, welche sich daselbst finden, anzuzeigen und zu verbessern. Durch Zerlegung von $\frac1{(x-1)(x+c)}$ in Partialbrüche und Anwendung der Formeln dieses Artikels findet man leicht
\[
\int_0^\infty\frac{x^a-1}{(x-1)(x+c)}\dd x
=\frac{\pi}{c+1}
\left(\frac{c^a-\cos a\pi}{\sin a\pi}
-\frac{\ell c}{\pi}\right),
\]
und diese Gleichung gilt für positive und negative echt gebrochene $a$, wovon man sich leicht überzeugt, wenn man $\frac1x$ und $\frac1c$ statt $x$ und $c$ schreibt. Differentiiert man in bezug auf $a$ und setzt dann $a=0$, so erhält man eine Reihe von Integralen, welche in jenen Tafeln unrichtig angegeben sind. Um die lästigen Differentiationen möglichst zu erleichtern, kann man folgenden Kunstgriff anwenden. Setzt man
\[
\frac{\ell c}{\pi}+\frac{c+1}{\pi}
\int_0^\infty\frac{x^a-1}{(x-1)(x+c)}\dd x
=\frac{c^a-\cos a\pi}{\sin a\pi}=J,
\qquad
\frac{\dd^nJ}{\dd a^n}={}^{n}J,
\]
so erhält man für $a=0$
\[
\int_0^\infty\frac{(\ell x)^n\dd x}{(x-1)(x+c)}
=\frac{\pi}{c+1}\,{}^{n}J_0.
\]
Man findet aber leicht
\[
\begin{aligned}
&J\pi^n\sin\left(a+\frac n2\right)\pi
+n\,{}^{1}J\,\pi^{n-1}\sin\left(a+\frac{n-1}{2}\right)\pi+\cdots\\
&\quad{}+n\,{}^{n-1}J\,\pi\sin\left(a+\frac12\right)\pi
+{}^{n}J\sin a\pi\\
&=\frac{\dd^n(J\sin a\pi)}{\dd a^n}
=c^a(\ell c)^n-\pi^n\cos\left(a+\frac n2\right)\pi,
\end{aligned}
\]
und für $a=0$ erhält man Rekursionsformeln für die $J_0$ mit geraden und die mit ungeraden Indizes. So findet man
\[
\begin{aligned}
\int_0^\infty\frac{\ell x\,\dd x}{(x-1)(x+c)}
&=\frac{(\ell c)^2+\pi^2}{2(c+1)},\\
\int_0^\infty\frac{(\ell x)^2\dd x}{(x-1)(x+c)}
&=\frac{\ell c\bigl((\ell c)^2+\pi^2\bigr)}{3(c+1)},\\
\int_0^\infty\frac{(\ell x)^3\dd x}{(x-1)(x+c)}
&=\frac{\bigl((\ell c)^2+\pi^2\bigr)^2}{4(c+1)},\\
\int_0^\infty\frac{(\ell x)^4\dd x}{(x-1)(x+c)}
&=\frac{\ell c\bigl((\ell c)^2+\pi^2\bigr)
        \bigl(3(\ell c)^2+7\pi^2\bigr)}{5(c+1)\cdot3},\\
\int_0^\infty\frac{(\ell x)^5\dd x}{(x-1)(x+c)}
&=\frac{\bigl((\ell c)^2+\pi^2\bigr)^2
        \bigl((\ell c)^2+3\pi^2\bigr)}{6(c+1)},
\end{aligned}
\]
während in jenen Tafeln statt der Divisoren $2,3,4,5,6$ die Divisoren $1\cdot2$, $1\cdot2\cdot3$, $1\cdot2\cdot3\cdot4$, $1\cdot2\cdot3\cdot4\cdot5$, $1\cdot2\cdot3\cdot4\cdot5\cdot6$ angegeben sind.

\end{document}
